Our pilot model lived inside a deliberately small corner of Python --- a
bit of arithmetic, a couple of comparisons, a few list operations, some
basic if/else. That tiny corner is what let us measure capability
cleanly: change one small thing at a time and watch what happens.

But a capability map drawn on one tiny corner only tells us where the
boundary is \emph{for that corner.} It would be a narrow conclusion if
we claimed ``this model can do \texttt{a\ +\ b} and nothing else'' ---
of course it could, we only taught it that. So the next question is the
natural one: \textbf{what happens when we widen the world?} Does giving
the model a richer language unlock more capability, or does it just give
it more ways to tie itself to surface form?

This is the chapter about \textbf{expanding the capability spec} --- the
versioned whitelist from Chapter 4 that defines what Python the model is
allowed to write.

\section{Why the world is versioned}\label{why-the-world-is-versioned}

The key mechanical fact we need here is that the capability spec is
\textbf{versioned} (\texttt{python\_capability\_v0\_1}, \texttt{v0\_2},
and so on). It is not a living, unwritten list of ``what the model can
do.'' It is a recorded, frozen definition of ``what the model is allowed
to do'' at a particular point.

This matters enormously for the one-variable discipline. If we widen the
allowed Python \emph{and} change the data \emph{and} retrain, we could
not say which change caused any shift in behavior. But if we record the
spec version alongside every result, we can point to exactly what grew
--- ``the model was given list comprehensions at spec v0\_3'' --- and
attribute any change to that. Versioning the world is what lets us make
a causal claim about capability, rather than a correlational guess. It
is the same reason we record the exact model configuration and dataset
in every manifest.

\section{How we widen the world}\label{how-we-widen-the-world}

Expanding the world is a controlled act, not a free-for-all. We do not
hand the model the entire Python standard library at once. We add
constructs deliberately, one family at a time, and always with a reason.
For example, we might add:

\begin{itemize}
\tightlist
\item
  a new built-in (say, a mathematical function we were keeping out),
\item
  a new control-flow construct (e.g.~\texttt{break}/\texttt{continue} in
  loops, if they were missing),
\item
  a new data structure (e.g.~dictionaries, if we had been lists-only),
\item
  short-circuit boolean logic, or a new comparison form.
\end{itemize}

Each addition is a \emph{step up a ladder}, not a teardown of the wall.
And each one changes the set of problems the model can even express. The
pilot's corner was deliberately small partly so that the \emph{boundary}
would be legible --- and partly so that, when we widen it later, we can
see the effect of \emph{just that} widening.

\section{A crucial reframing: syntax is not the hard
part}\label{a-crucial-reframing-syntax-is-not-the-hard-part}

Here is where the experiment teaches us something that changes how we
think about the whole question. It is tempting to imagine that the
model's limit is about \emph{syntax} --- that if we just teach it more
Python, it will be more capable. Our results say otherwise, and it
matters.

Consider the composition result from Chapter 16. The model knew the
syntax of both \texttt{abs} and \texttt{add}. It could produce
\texttt{return\ a\ +\ b} and \texttt{return\ abs(x)} perfectly. The
break came when it needed to \emph{combine} them into
\texttt{return\ abs(a\ +\ b)}. There was no new syntax involved ---
every token in \texttt{abs(a\ +\ b)} was already within the model's
known vocabulary and its allowed subset. The model simply could not put
two things it knew together.

So the real dividing line in this book is \textbf{not ``how much Python
the model is given.''} It is \textbf{``whether the model can compose the
Python it already knows.''} Giving the model \emph{more} syntax would
let it express more kinds of programs, but it would not give it the
\emph{ability to combine} --- and that ability, not the vocabulary, is
what separates a bag of reflexes from genuine capability.

\section{Why we still widen the
world}\label{why-we-still-widen-the-world}

That said, widening the world remains the right thing to do, and for an
honest reason: it lets us \emph{test whether the boundary moves}. The
interesting scientific question is not ``can a tiny model learn any
fixed corner?'' --- the answer is yes. The question is ``does the
boundary of what it can do shift when we push the world's edges out?''
If we widen the allowed Python and the model still cannot compose, that
is stronger evidence that the limit is structural, not about the surface
of the language. If the boundary does move, that is a genuine discovery.

So this chapter is less about a single dramatic result and more about a
discipline: a controlled, versioned, one-variable-at-a-time
\emph{capacity ladder} that lets us watch exactly where the model starts
to fall apart as we widen what it may do.

\section{The ladder, as a discipline}\label{the-ladder-as-a-discipline}

The discipline is: define one more capability, generate data for it,
train, evaluate, audit, record --- \emph{then} move on. Never widen
multiple things at once. The ladder is not a claim that the model will
reach the top; it is a map of where the model \emph{does} stop.

And a stopping point is itself a result. We found ours at composition:
the model can do each rung's primitives but not combine them. That
finding --- cheap and specific --- is worth more than a glossy claim of
having ``reached Python Level 8.'' Knowing exactly where the model
stops, and being able to say \emph{why}, is the entire value of this
kind of careful, pessimistic, honest investigation.

\section{The honest summary}\label{the-honest-summary}

Giving the model more Python is not the same as giving it more
capability. The pilot already revealed the real boundary: composition.
Widening the world is worth doing --- but the purpose is to \emph{test
where the boundary is and whether it moves}, not to assume that more
syntax automatically means more understanding.
